\documentclass[10pt,a4paper]{amsart}
\usepackage[margin=19mm]{geometry}
\usepackage{amsmath,amssymb,mathtools}
\usepackage[scr=rsfs,cal=euler]{mathalfa}
\usepackage{xfrac}
\usepackage[unicode,hidelinks,bookmarks=false]{hyperref}
\newcommand{\ie}{i.e.\ }
\newcommand{\cf}{cf.\ }
\newcommand{\resp}{respectively\ }
\newcommand{\Subsection}{\subsection}
\providecommand{\nobreakdash}{\nobreak\hbox{-}}
\setlength{\emergencystretch}{2em}
\multlinegap0pt
\theoremstyle{plain}
\newtheorem{theorem}{Theorem}[section]
\newtheorem{lemma}[theorem]{Lemma}
\newtheorem{proposition}[theorem]{Proposition}
\newtheorem{corollary}[theorem]{Corollary}
\theoremstyle{definition}
\newtheorem{definition}[theorem]{Definition}
\theoremstyle{remark}
\newtheorem{remark}[theorem]{Remark}
\title{When Concealed Links Cannot Be Recovered: A Structural Identifiability Bound and Evaluation Pitfalls in Offshore Leak Networks}
\author{Joseph Bingham}
\date{Source: arXiv:2609.26171v1 (CC BY 4.0); the author's own native \LaTeX{} source, set here in the editorial AMS \texttt{amsart} wrapper (the publisher's ACM \texttt{acmart} class named by the source is not redistributed).}
\begin{document}
\maketitle
\noindent\textit{Editorial scope.} The counted claim of this unit is the article's Theorem 4.1 (source label \texttt{thm:floor}), the indistinguishability floor for structural scoring rules, delivered together with the author's complete original proof. The statement and the proof are the only retained passages and every one of their characters is the author's own; the proof is delivered as the single primary proof body exactly as the author's proof environment encloses it. Printed numbering: the publisher's class declares the theorem environment, and that class is not present in this package and is not redistributed, so the wrapper restates the declaration, and the section and theorem counters are advanced so that the delivered PDF prints the counted claim as Theorem~4.1, exactly as the published article does. Omitted: the rest of the article, that is its title block, abstract, introduction, background section, the data and problem-formulation section that precedes the counted theorem and everything that follows it. The retained unit is the maximal source-local passage of the article that is closed under its own cross-references: the surrounding text refers to notation and to figures of the article that cannot be retained inside a unit of this size, and the counted statement and proof are self-contained. Nothing of the counted statement or of the counted proof is omitted. Class substitution: the source class is the publisher's ACM \texttt{acmart} class in sigconf mode; neither that class nor any of its style files is present in this package and none is redistributed, no publisher logo or artwork is redistributed, and the retained author text is set in an editorial AMS \texttt{amsart} wrapper with the author's own mathematical notation. Reference handling: no retained passage contains a citation; the article's bibliography is part of the author's own TeX and is bundled unchanged inside the source file, and no reference entry is transcribed, delivered or invented. No mathematics is written, repaired or re-derived here.

\setcounter{section}{4}
\setcounter{theorem}{0}
\begin{theorem}[Indistinguishability floor]\label{thm:floor}
Fix a target entity $e$, let $H=\mathrm{Stab}(e)$, and let $C=H\!\cdot\!o$ be the
orbit of the true owner $o$ under $H$. For every structural rule $s$, the score
$s(\cdot,e;G)$ is constant on $C$. Two consequences follow.
\begin{enumerate}
\item[(i)] Ranking $o$ against a uniformly random member of $C$ gives expected
ROC-\textsc{auc} exactly $\tfrac12$.
\item[(ii)] The probability that $o$ is ranked strictly first among all
candidates is at most $1/|C|$.
\end{enumerate}
\end{theorem}

\begin{proof}
Take any $o',o''\in C$. By definition of the orbit there is $\pi\in H$ with
$\pi(o')=o''$, and because $\pi$ fixes $e$ we have $\pi(e)=e$. Isomorphism
invariance then gives $s(o',e;G)=s(\pi(o'),\pi(e);G)=s(o'',e;G)$, so $s(\cdot,e;G)$
takes a single value across $C$. A constant score orders the members of $C$
uniformly at random. For (i), the expected \textsc{auc} of one distinguished
member against a random other under a uniform order is $\tfrac12$. For (ii), the
true owner is equally likely to occupy any position among the $|C|$ tied
candidates, so it lands first with probability at most $1/|C|$.
\end{proof}

\end{document}
